% FreeCAD macOS-Native + CAD-Familiar UI - Scoping Document
% Minimalist, Apple-inspired styling. Compile with: pdflatex macos-native-ui-scope.tex
\documentclass[11pt]{article}

% --- Page geometry: generous whitespace -------------------------------------
\usepackage[a4paper,top=26mm,bottom=24mm,left=24mm,right=24mm]{geometry}

% --- Typography: clean sans-serif throughout --------------------------------
\usepackage[T1]{fontenc}
\usepackage[utf8]{inputenc}
\usepackage[scaled=1.02]{helvet}
\usepackage[scaled=0.92]{inconsolata}  % narrow, clean monospace for paths
\renewcommand{\familydefault}{\sfdefault}
\usepackage{microtype}
\usepackage[htt]{hyphenat}             % allow breaking inside typewriter text
\usepackage[document]{ragged2e}
\usepackage{setspace}
\setstretch{1.18}
% Let long file paths break gracefully rather than overrun the right margin
\sloppy
\emergencystretch=3em

% --- Muted, restrained color palette ----------------------------------------
\usepackage{xcolor}
\definecolor{appleink}{HTML}{1D1D1F}
\definecolor{applegray}{HTML}{6E6E73}
\definecolor{appleblue}{HTML}{0071E3}
\definecolor{applegreen}{HTML}{34C759}
\definecolor{appleamber}{HTML}{FF9F0A}
\definecolor{applered}{HTML}{FF3B30}
\definecolor{hairline}{HTML}{D2D2D7}
\definecolor{panel}{HTML}{F5F5F7}
\color{appleink}

% --- Section styling: light, spacious ---------------------------------------
\usepackage{titlesec}
\setcounter{secnumdepth}{0}
\titleformat{\section}{\LARGE\bfseries\color{appleink}}{}{0em}{}
\titlespacing*{\section}{0pt}{2.4em}{0.8em}
\titleformat{\subsection}{\large\bfseries\color{appleink}}{}{0em}{}
\titlespacing*{\subsection}{0pt}{1.4em}{0.4em}
\titleformat{\subsubsection}{\normalsize\bfseries\color{applegray}}{}{0em}{}
\titlespacing*{\subsubsection}{0pt}{1.0em}{0.2em}

% --- Lists ------------------------------------------------------------------
\usepackage{enumitem}
\setlist{leftmargin=1.3em,itemsep=3pt,topsep=4pt,parsep=0pt}
\setlist[itemize]{label=\textcolor{appleblue}{\small\raisebox{0.15ex}{\textbullet}}}

% --- Tables -----------------------------------------------------------------
\usepackage{booktabs}
\usepackage{tabularx}
\usepackage{array}
\renewcommand{\arraystretch}{1.3}
\newcolumntype{L}{>{\RaggedRight\arraybackslash}X}

% --- Boxes ------------------------------------------------------------------
\usepackage[most]{tcolorbox}

% --- Links ------------------------------------------------------------------
\usepackage[hidelinks]{hyperref}

% --- Headers / footers ------------------------------------------------------
\usepackage{fancyhdr}
\pagestyle{fancy}
\fancyhf{}
\renewcommand{\headrulewidth}{0pt}
\renewcommand{\footrulewidth}{0pt}
\fancyfoot[C]{\footnotesize\color{applegray}FreeCAD\quad\textbullet\quad macOS-Native \& CAD-Familiar UI\quad\textbullet\quad\thepage}

\newcommand{\hairrule}{\vspace{2pt}{\color{hairline}\hrule height 0.6pt}\vspace{2pt}}
\usepackage{soul}
\newcommand{\eyebrow}[1]{{\footnotesize\color{appleblue}\bfseries\so{\MakeUppercase{#1}}}}

% Traffic-light dots for the title flourish
\newcommand{\trafficlights}{%
  \raisebox{0.2ex}{\textcolor{applered}{\Large\textbullet}}\,%
  \raisebox{0.2ex}{\textcolor{appleamber}{\Large\textbullet}}\,%
  \raisebox{0.2ex}{\textcolor{applegreen}{\Large\textbullet}}%
}

% Priority tag
\newtcbox{\ptag}[1][appleblue]{on line, colframe=#1, colback=#1, coltext=white,
  boxrule=0pt, arc=3pt, left=4pt, right=4pt, top=1pt, bottom=1pt,
  fontupper=\footnotesize\bfseries}

\begin{document}
\thispagestyle{fancy}

% ============================ TITLE BLOCK ===================================
\begingroup
\setstretch{1.0}
\noindent\trafficlights\hfill\eyebrow{Product Scoping}\\[1.0em]
{\fontsize{29}{33}\selectfont\bfseries macOS-Native Experience \&\\ CAD-Familiar Defaults\par}
\vspace{0.7em}
{\large\color{applegray} Making FreeCAD feel at home on Apple Silicon, and instantly familiar to SolidWorks and Fusion 360 users.\par}
\vspace{1.1em}
\noindent{\footnotesize\color{applegray} Repository \; \textbf{brucecodes2023/FreeCAD\_BT} \hfill Platform focus \; \textbf{macOS (arm64)} \hfill Status \; \textbf{Scoping}}
\endgroup

\vspace{0.9em}
\hairrule
\vspace{1.3em}

% ============================ SUMMARY PANEL =================================
\begin{tcolorbox}[enhanced, colback=panel, colframe=panel, boxrule=0pt,
  arc=6pt, left=16pt, right=16pt, top=12pt, bottom=12pt]
{\eyebrow{Summary}}\\[0.6em]
This document scopes three workstreams, ordered by the priority you set: (A) make
FreeCAD's \emph{default} UI feel like a modern mechanical CAD tool in the
SolidWorks / Fusion 360 tradition; (B) sharpen the \emph{macOS-native} window
experience; and (C) make FreeCAD \emph{faster on Apple Silicon}. Each item below is
grounded in the current codebase, notes what already exists versus what is new
work, lists the files it touches, and rates effort by technical scope (not
calendar time).

\vspace{0.6em}
\textbf{Two headline findings shape the plan:}
\begin{itemize}
  \item A \textbf{SolidWorks navigation style already exists} in-tree, so a familiar
  mouse/orbit model is mostly a \emph{defaults} change, not new code.
  \item FreeCAD's main window \textbf{already uses the native macOS frame}, so the
  red/amber/green traffic lights are \emph{already} rendered top-left by macOS. The
  real ``native'' work is unifying the toolbar with the title bar and de-styling
  non-native chrome \textemdash{} not moving the buttons.
\end{itemize}
\end{tcolorbox}

% ============================ OBJECTIVES ====================================
\section{Objectives \& Priority}
\begin{enumerate}[label=\textbf{\arabic*.}]
  \item \textbf{CAD-familiar defaults first.} Out of the box, navigation, view cube,
  layout, and theme should read as a modern mechanical-CAD tool to a SolidWorks /
  Fusion user.
  \item \textbf{Native macOS window chrome.} Title bar, toolbar, and panels should
  follow macOS conventions; verify the traffic-light placement is already correct.
  \item \textbf{Apple Silicon performance.} Native arm64 build tuning, smarter
  startup, and a path toward a modern GPU backend.
\end{enumerate}

% ============================ WORKSTREAM A =================================
\section{Workstream A \; \ptag{Priority 1}\; CAD-Familiar UI Defaults}

Goal: a first launch that feels like SolidWorks / Fusion 360 without the user
touching a single preference. The cleanest, lowest-risk vehicle is a built-in
\textbf{Preference Pack} plus a small change to the shipped defaults.

\subsection{A1 \textemdash{} Default to a familiar navigation model}
\textbf{Exists today.} Navigation styles live under \texttt{src/Gui/Navigation/}; the
current default is \texttt{CADNavigationStyle}. A \texttt{SolidWorksNavigationStyle} is
\emph{already implemented and registered} (via \texttt{src/Gui/SoFCDB.cpp}). The default
is read from the parameter \texttt{BaseApp/Preferences/View$\rightarrow$NavigationStyle},
with the code fallback set in \texttt{src/Gui/View3DSettings.cpp} and the template
default in \texttt{src/Gui/PreferencePackTemplates/View.cfg}.
\begin{itemize}
  \item \textbf{Change:} ship \texttt{NavigationStyle = Gui::SolidWorksNavigationStyle}
  as the default (via the new pack in A2, or the fallback in \texttt{View3DSettings.cpp}).
  \item \textbf{Effort:} Low. \quad\textbf{Risk:} Low (reversible preference).
\end{itemize}

\subsection{A2 \textemdash{} Ship a ``Modern CAD'' preference pack}
\textbf{New work, built on existing infrastructure.} Preference packs live under
\texttt{src/Gui/PreferencePacks/} and are managed by \texttt{Gui::PreferencePackManager}
(\texttt{src/Gui/PreferencePacks/PreferencePackManager.cpp}). A pack is a folder with a
\texttt{.cfg} (merged into user parameters), optional \texttt{pre}/\texttt{post} macros, and
optional \texttt{.qss} stylesheet. Bundle into one pack:
\begin{itemize}
  \item navigation style (A1) and view-cube configuration (A3);
  \item a toolbar/dock layout derived from \texttt{PreferencePackTemplates/Main\_window\_layout.cfg} (tree, properties, and a modeling-forward default);
  \item a light + dark stylesheet in the SolidWorks/Fusion visual register, alongside \texttt{src/Gui/Stylesheets/FreeCAD.qss}.
\end{itemize}
\textbf{Effort:} Medium. \quad\textbf{Risk:} Low\textendash Medium (layout/theme QA).

\subsection{A3 \textemdash{} View cube tuned like a design tool}
\textbf{Exists today.} The navigation cube (\texttt{src/Gui/NaviCube.cpp},
\texttt{src/Gui/Inventor/SoNaviCube.h}) is already a SolidWorks/Fusion-style view cube.
It is fully configurable via \texttt{BaseApp/Preferences/NaviCube} (corner, size, colors,
fonts). Set sensible defaults (corner, size, subtle colors) inside the A2 pack.
\textbf{Effort:} Low.

\subsection{A4 \textemdash{} First-run onboarding hook}
The Start workbench already has a first-run flow (\texttt{FirstStart2024} flag) with a
theme selector: \texttt{src/Mod/Start/Gui/GeneralSettingsWidget.cpp} and
\texttt{ThemeSelectorWidget}. Add the ``Modern CAD'' pack as a selectable (or default)
option so new users land in the familiar experience.
\textbf{Effort:} Low\textendash Medium.

\subsection{A5 \textemdash{} Stretch: a true Fusion 360 navigation style}
\textbf{New work.} No \texttt{Fusion360NavigationStyle} exists. Fusion's input model
(MMB-orbit, Shift+MMB-pan, scroll-zoom) differs from SolidWorks. Add a
\texttt{UserNavigationStyle} subclass under \texttt{src/Gui/Navigation/} and register it in
\texttt{src/Gui/SoFCDB.cpp} + \texttt{src/Gui/CMakeLists.txt}.
\textbf{Effort:} Medium. \quad\textbf{Risk:} Low (additive).

\subsection{A6 \textemdash{} Out of near-term scope: ribbon UI}
Neither SolidWorks' CommandManager nor Fusion's toolbar is present; FreeCAD uses a
classic \texttt{Gui::ToolBarManager} model (\texttt{src/Gui/ToolBarManager.cpp}) with a
\texttt{WorkbenchSelector}. A ribbon is greenfield and large; recommend deferring in
favor of A1\textendash A4, which capture most of the ``familiar'' feel at a fraction of
the risk.

% ============================ WORKSTREAM B =================================
\section{Workstream B \; \ptag[applegray]{Priority 2}\; macOS-Native Window Chrome}

\subsection{B0 \textemdash{} Finding: traffic lights are already native \& top-left}
\texttt{Gui::MainWindow} (\texttt{src/Gui/MainWindow.cpp}) is a plain \texttt{QMainWindow} with
default \texttt{Qt::Window} flags \textemdash{} no frameless or custom title bar, and no
\texttt{setUnifiedTitleAndToolBarOnMac} anywhere in the tree. On macOS, Qt/Cocoa draws
the standard frame, so the red/amber/green controls already sit top-left. If they
appear top-right, the app is being viewed on Linux/Windows or through a themed
build. \textbf{Action:} verify visually on macOS (Qt~6.8) and document; no code move
is required to place the buttons.

\subsection{B1 \textemdash{} Unified title + toolbar}
Add \texttt{setUnifiedTitleAndToolBarOnMac(true)} in the \texttt{MainWindow} constructor to
merge the top toolbar into the title bar, the hallmark of a native Mac app.
Validate toolbar layout, full-screen, and the \texttt{ToolBarManager} areas.
\textbf{Effort:} Low\textendash Medium. \quad\textbf{Risk:} Medium (layout regressions).

\subsection{B2 \textemdash{} Native dock \& panel chrome}
Dock panels currently draw custom title bars via \texttt{QDockWidget::title} in
\texttt{src/Gui/Stylesheets/FreeCAD.qss}, and floating toolbars use
\texttt{Qt::FramelessWindowHint} (\texttt{src/Gui/ToolBarManager.cpp}). Platform-gate this
QSS on macOS so Cocoa draws native dock/panel chrome.
\textbf{Effort:} Medium. \quad\textbf{Risk:} Medium (theme interplay).

\subsection{B3 \textemdash{} Full-size content view (deep native integration)}
For a truly modern Mac look (content under the title bar, large-title behavior),
add \texttt{NSWindow} integration via a new Objective-C++ (\texttt{.mm}) shim. There is a
precedent for macOS Obj-C in-tree (QuickLook under \texttt{src/MacAppBundle/QuickLook/},
3Dconnexion \texttt{GuiNativeEventMac}), but \emph{no} window-chrome Cocoa code exists yet.
\textbf{Effort:} High. \quad\textbf{Risk:} High (new native surface).

\subsection{B4 \textemdash{} Bundle identity \& appearance}
The bundled \texttt{src/MacAppBundle/FreeCAD.app/Contents/Info.plist} ships an empty
bundle identifier (filled by CMake \texttt{sed}); the conda recipe uses
\texttt{org.freecad.FreeCAD} and sets \texttt{NSRequiresAquaSystemAppearance=False}
(\texttt{package/rattler-build/osx/Info.plist.template}). Align a stable bundle id,
dark-mode support, and icon (\texttt{freecad.icns}); track signing/notarization as a
release concern.
\textbf{Effort:} Low\textendash Medium.

% ============================ WORKSTREAM C =================================
\section{Workstream C \; \ptag[applegray]{Priority 3}\; Apple Silicon Performance}

\subsubsection{What already helps}
Native arm64 comes from conda-forge/\texttt{pixi} (\texttt{osx-arm64} platform in
\texttt{pixi.toml}); \texttt{ccache} is on by default (\texttt{CMakeLists.txt}); C++ workbench
modules load lazily on activation; and hot paths (mesh, TechDraw HLR) already use
\texttt{QtConcurrent}.

\subsection{C1 \textemdash{} Pin architecture \& deployment target}
No \texttt{CMAKE\_OSX\_ARCHITECTURES} is set; the build follows the toolchain. Set
\texttt{arm64} explicitly for reproducible native builds and confirm
\texttt{CMAKE\_OSX\_DEPLOYMENT\_TARGET} (package build defaults to 11.0). Files:
\texttt{CMakePresets.json} (\texttt{conda-macos-*}), \texttt{package/rattler-build/build.sh}.
\textbf{Effort:} Low.

\subsection{C2 \textemdash{} Optimization: LTO/IPO \& Apple CPU tuning}
No LTO/IPO or \texttt{-mcpu=apple-m*} tuning in the main CMake. Add optional
\texttt{INTERPROCEDURAL\_OPTIMIZATION} and Apple-CPU flags for release builds
(\texttt{cMake/FreeCAD\_Helpers/SetGlobalCompilerAndLinkerSettings.cmake}). Measure
binary size and link time.
\textbf{Effort:} Medium. \quad\textbf{Risk:} Medium (build-time/ABI QA).

\subsection{C3 \textemdash{} Activate TBB parallelism}
\texttt{tbb-devel} is already a dependency but \texttt{HAVE\_TBB} is never defined, so the
TBB code paths (e.g. \texttt{src/Mod/Import/App/ImportOCAF.cpp}) are dead. Wire
\texttt{find\_package(TBB)} + \texttt{HAVE\_TBB} in \texttt{src/config.h.cmake} to light up
parallel import/mesh paths on multi-core Apple Silicon.
\textbf{Effort:} Medium. \quad\textbf{Risk:} Medium.

\subsection{C4 \textemdash{} Faster startup}
All ${\sim}28$ modules' \texttt{InitGui.py} scripts run eagerly at GUI startup
(\texttt{src/Gui/FreeCADGuiInit.py}). Explore deferring non-essential workbench
registration and leaning on the existing \texttt{BackgroundAutoloadModules} mechanism
(\texttt{src/Gui/StartupProcess.cpp}) to cut time-to-interactive.
\textbf{Effort:} Medium. \quad\textbf{Risk:} Medium (ordering assumptions).

\subsection{C5 \textemdash{} Long-term: modern GPU backend}
Rendering is OpenGL-only through Coin3D + \texttt{QOpenGLWidget}
(\texttt{src/Gui/Quarter/QuarterWidget.cpp}, \texttt{src/Gui/View3DInventorViewer.cpp});
macOS OpenGL is deprecated. A Metal path (via Qt RHI or MoltenVK) is a substantial
Coin + Quarter + viewer effort, not a toggle. Flag as strategic R\&D, decoupled
from A/B.
\textbf{Effort:} Very High. \quad\textbf{Risk:} High.

% ============================ SEQUENCING ===================================
\section{Sequencing}
\noindent\begin{tabularx}{\linewidth}{@{}p{0.14\linewidth}p{0.30\linewidth}L@{}}
\toprule
\textbf{Phase} & \textbf{Focus} & \textbf{Items} \\
\midrule
Phase 1 & Familiar defaults & A1, A3, A2 (pack), A4 onboarding \\
Phase 2 & Native window feel & B0 verify, B1 unified toolbar, B4 bundle id \\
Phase 3 & Native polish + speed & B2 dock chrome, C1, C3, C4 \\
Phase 4 & Strategic R\&D & A5 Fusion nav, B3 NSWindow, C5 Metal \\
\bottomrule
\end{tabularx}

% ============================ VALIDATION ===================================
\section{Validation Plan}
\noindent\begin{tabularx}{\linewidth}{@{}p{0.30\linewidth}L@{}}
\toprule
\textbf{Change} & \textbf{How it is proven} \\
\midrule
Default navigation / pack & Fresh profile launches with SolidWorks-style orbit and the Modern CAD layout; toggling back works. \\
View cube defaults & Cube appears in the configured corner with the intended size/colors. \\
Traffic lights (B0) & Screenshot on macOS showing native top-left controls; documented. \\
Unified toolbar (B1) & Visual check of title/toolbar merge; full-screen and toolbar-area regression pass. \\
arm64 / LTO / TBB & Build reports target \texttt{arm64}; benchmark import/mesh with and without TBB; record link time. \\
Startup (C4) & Time-to-interactive measured before/after, across representative workbenches. \\
No regressions & FreeCAD CLI + C++ self-tests remain green (as in the environment validation). \\
\bottomrule
\end{tabularx}

% ============================ FILE INDEX ===================================
\section{Grounded File Index}
\noindent\begin{tabularx}{\linewidth}{@{}p{0.42\linewidth}L@{}}
\toprule
\textbf{Concern} & \textbf{Key path(s)} \\
\midrule
Navigation styles (SolidWorks exists) & \texttt{src/Gui/Navigation/*}, \texttt{src/Gui/SoFCDB.cpp} \\
Navigation default \& fallback & \texttt{src/Gui/View3DSettings.cpp}, \texttt{PreferencePackTemplates/View.cfg} \\
View cube & \texttt{src/Gui/NaviCube.cpp}, \texttt{src/Gui/Inventor/SoNaviCube.h} \\
Preference packs \& manager & \texttt{src/Gui/PreferencePacks/}, \texttt{PreferencePackManager.cpp} \\
Stylesheets / theme & \texttt{src/Gui/Stylesheets/FreeCAD.qss}, \texttt{defaults.qss} \\
First-run onboarding & \texttt{src/Mod/Start/Gui/GeneralSettingsWidget.cpp}, \texttt{ThemeSelectorWidget} \\
Main window \& chrome & \texttt{src/Gui/MainWindow.cpp}, \texttt{src/Gui/ToolBarManager.cpp} \\
macOS bundle & \texttt{src/MacAppBundle/}, \texttt{package/rattler-build/osx/} \\
Build presets / arch & \texttt{CMakePresets.json}, \texttt{pixi.toml}, \texttt{package/rattler-build/build.sh} \\
Compiler / linker flags & \texttt{cMake/FreeCAD\_Helpers/SetGlobalCompilerAndLinkerSettings.cmake} \\
TBB gate & \texttt{src/config.h.cmake}, \texttt{src/Mod/Import/App/ImportOCAF.cpp} \\
Rendering stack & \texttt{src/Gui/Quarter/QuarterWidget.cpp}, \texttt{src/Gui/View3DInventorViewer.cpp} \\
GUI startup & \texttt{src/Gui/FreeCADGuiInit.py}, \texttt{src/Gui/StartupProcess.cpp} \\
\bottomrule
\end{tabularx}

\vspace{1.2em}
\hairrule
\vspace{0.6em}
\noindent{\footnotesize\color{applegray} Effort ratings describe technical scope (breadth of components and invasiveness), not calendar estimates. ``Exists today'' items are configuration/defaults changes; ``new work'' items add code behind existing extension points.}

\end{document}
